\chapter{Nơ-ron không phải một thiết bị duy nhất}
% full version: chapters/18_eetypes.tex

Ở chương một, chúng ta đã phân loại nơ-ron theo chất mà chúng nhả ra. Một kỹ sư sẽ phân loại khác: mỗi nơ-ron hoạt động như thiết bị nào. Và hóa ra dưới vẻ ngoài giống nhau lại ẩn cả một hộp linh kiện điện tử.

\section*{Hộp linh kiện điện tử}

Cái xô thủng ở chương hai là một bộ tích lũy có rò: nó cộng các tín hiệu trong một cửa sổ ngắn. Với cửa sổ rất nhỏ, nó trở thành bộ phát hiện trùng khớp. Nơ-ron chỉ đáp ứng với thay đổi rồi im ngay là một bộ lọc cho sự thay đổi đi qua. Nơ-ron tự hoạt động là máy đếm nhịp, một bộ dao động; nếu nhịp của nó được một đầu vào vặn, đó là bộ dao động điều khiển bằng điện áp. Nơ-ron phát ra chùm tiếng tách là bộ tạo chùm xung, viết ``dấu gạch''. Nơ-ron tách theo nhịp sóng âm là bộ tách pha.

\pencilfig{18}{\pencilart{18}}{Dưới vẻ ngoài giống nhau ẩn cả một hộp linh kiện điện tử: điện trở, tụ điện, cuộn cảm, bộ dao động.}

Còn có những thiết bị mà trước đây chúng ta chưa nói tới. \emph{Bộ cộng hưởng} giống một chiếc xích đu: nó đáp ứng tốt nhất với những cú đẩy theo một nhịp nhất định, như chiếc xích đu chỉ đưa cao khi được đẩy đúng lúc. Bên trong nơ-ron, vai trò của lò xo do một dòng điện chậm chống lại sự thay đổi đảm nhận. \emph{Bộ tích lũy không rò} giữ giá trị lâu, như một tụ điện không rò --- đó là cách hoạt động của mạng giữ vị trí mắt khi bạn nhìn sang bên. Nhưng muốn vậy, phản hồi phải được chỉnh gần như hoàn hảo, và đó là một sự tinh chỉnh tốn kém. \emph{Bộ chốt} là nơ-ron mà một cú đẩy ngắn chuyển sang hoạt động dài, cho tới khi bị tắt, như một công tắc có khóa giữ; ở các nơ-ron điều khiển cơ, chế độ này được serotonin bật lên.

\section*{Thiết bị tắc kè hoa}

Phát hiện chính của cách phân loại này: đây không phải những linh kiện khác nhau mà là những \emph{chế độ} khác nhau của cùng một linh kiện. Cùng một nơ-ron có thể là bộ tích lũy ban ngày và bộ tạo chùm xung ban đêm, tùy theo điều các đài phát thanh nói với nó. Một kỹ sư sẽ gọi bộ não là một ma trận logic lập trình lại được: mạch chỉ có một, còn chức năng của linh kiện do các thiết lập thay đổi.

Việc nơ-ron làm việc được ở tất cả các chế độ này đã được đo. Còn bộ não dùng việc chuyển chế độ để tính toán ra sao thì phần lớn vẫn là giả thuyết.

\fullversion{18}
